% GENERATED FILE. Edit canonical lessons, then run npm run generate:books.
% canonical-source-hash: fnv1a64:4fbc1a42ea32416e
% canonical-lessons: BN-C47-pishi, BN-C47-shbami, BN-C47-stri, BN-C47-shishu, BN-C47-lok

\chapter{Aunt, Husband, Wife, Child, Person}
\label{ch:bn-aunt-husband-wife-child-person}
\hlchaptermodality{47}

\begin{chapteropening}
\textbf{By the end of this chapter:} \emph{I can say an aunt, a husband, a wife, a child and a person.}

Complete the last lesson of chapter 47: I can say an aunt, a husband, a wife, a child and a person.
\end{chapteropening}

\section[\texorpdfstring{\bn{পিসি} (pishi)}{পিসি (pishi)}]{\bn{পিসি} (pishi) — an aunt (father's sister)}

\label{lesson:BN-C47-pishi}



\begin{quote}
Before the new one: say the Bengali for an uncle, then the Bengali for an aunt.
\end{quote}

\subsection*{You'll want to know: \bn{পিসি}}
\textbf{\bn{পিসি}} — \emph{pishi} — \textquotedblleft{}an aunt (father's sister)\textquotedblright{}.

\begin{cousinweb}
\textbf{\bn{পিসি}} (\emph{pishi}) is a tadbhava, a word inherited from Sanskrit and reshaped by sound change. It comes from \emph{pituḥṣvasṛ}, \textquotedblleft{}father's sister\textquotedblright{} (CDIAL 8177). It is \emph{pitṛ}, \textquotedblleft{}father\textquotedblright{}, + \emph{svasṛ}, \textquotedblleft{}sister\textquotedblright{}.
\end{cousinweb}

\begin{grammarlens}[title={What you've built}]
One more word for this chapter.
\end{grammarlens}

\subsection*{Your turn}
\begin{itemize}
\raggedright
  \item \emph{Say it:} \emph{pishi}
  \item \emph{Say it:} \emph{pishi}, once more
  \item \emph{Say it:} say \emph{pishi}
  \item \emph{Recall:} say the Bengali for an uncle, then the Bengali for an aunt, then say \emph{pishi} again
\end{itemize}

\begin{tcolorbox}[breakable,colback=teal!4,colframe=teal!35!black,title={Before you move on}]
What does \bn{পিসি} mean? (\textquotedblleft{}An aunt (father's sister)\textquotedblright{}.) Say it once more.
\end{tcolorbox}
\section[\texorpdfstring{\bn{স্বামী} (shbāmi)}{স্বামী (shbāmi)}]{\bn{স্বামী} (shbāmi) — a husband}

\label{lesson:BN-C47-shbami}



\begin{quote}
Before the new one: say the Bengali for an aunt, then the Bengali for an aunt.
\end{quote}

\subsection*{You'll want to know: \bn{স্বামী}}
\textbf{\bn{স্বামী}} — \emph{shbāmi} — \textquotedblleft{}a husband\textquotedblright{}.

\begin{cousinweb}
\textbf{\bn{স্বামী}} (\emph{shbāmi}) is Sanskrit \emph{svāmin}, \textquotedblleft{}owner, master; husband\textquotedblright{}, taken over whole as a learned loan, a tatsama (CDIAL 13930). The word inherited from it is \emph{sāṁi}, \textquotedblleft{}master\textquotedblright{}, which CDIAL lists.
\end{cousinweb}

\begin{grammarlens}[title={What you've built}]
One more word for this chapter.
\end{grammarlens}

\subsection*{Your turn}
\begin{itemize}
\raggedright
  \item \emph{Say it:} \emph{shbāmi}
  \item \emph{Say it:} \emph{shbāmi}, once more
  \item \emph{Say it:} say \emph{shbāmi}
  \item \emph{Recall:} say the Bengali for an aunt, then the Bengali for an aunt, then say \emph{shbāmi} again
\end{itemize}

\begin{tcolorbox}[breakable,colback=teal!4,colframe=teal!35!black,title={Before you move on}]
What does \bn{স্বামী} mean? (\textquotedblleft{}A husband\textquotedblright{}.) Say it once more.
\end{tcolorbox}
\section[\texorpdfstring{\bn{স্ত্রী} (stri)}{স্ত্রী (stri)}]{\bn{স্ত্রী} (stri) — a wife}

\label{lesson:BN-C47-stri}



\begin{quote}
Before the new one: say the Bengali for an aunt, then the Bengali for a husband.
\end{quote}

\subsection*{You'll want to know: \bn{স্ত্রী}}
\textbf{\bn{স্ত্রী}} — \emph{stri} — \textquotedblleft{}a wife\textquotedblright{}.

\begin{cousinweb}
\textbf{\bn{স্ত্রী}} (\emph{stri}) is Sanskrit \emph{strī}, \textquotedblleft{}woman, wife\textquotedblright{}, taken over whole as a learned loan, a tatsama.
\end{cousinweb}

\begin{grammarlens}[title={What you've built}]
One more word for this chapter.
\end{grammarlens}

\subsection*{Your turn}
\begin{itemize}
\raggedright
  \item \emph{Say it:} \emph{stri}
  \item \emph{Say it:} \emph{stri}, once more
  \item \emph{Say it:} say \emph{stri}
  \item \emph{Recall:} say the Bengali for an aunt, then the Bengali for a husband, then say \emph{stri} again
\end{itemize}

\begin{tcolorbox}[breakable,colback=teal!4,colframe=teal!35!black,title={Before you move on}]
What does \bn{স্ত্রী} mean? (\textquotedblleft{}A wife\textquotedblright{}.) Say it once more.
\end{tcolorbox}
\section[\texorpdfstring{\bn{শিশু} (shishu)}{শিশু (shishu)}]{\bn{শিশু} (shishu) — a child}

\label{lesson:BN-C47-shishu}



\begin{quote}
Before the new one: say the Bengali for a husband, then the Bengali for a wife.
\end{quote}

\subsection*{You'll want to know: \bn{শিশু}}
\textbf{\bn{শিশু}} — \emph{shishu} — \textquotedblleft{}a child\textquotedblright{}.

\begin{cousinweb}
\textbf{\bn{শিশু}} (\emph{shishu}) is Sanskrit \emph{śiśu}, \textquotedblleft{}young of an animal, child\textquotedblright{}, taken over whole as a learned loan, a tatsama.
\end{cousinweb}

\begin{grammarlens}[title={What you've built}]
One more word for this chapter.
\end{grammarlens}

\subsection*{Your turn}
\begin{itemize}
\raggedright
  \item \emph{Say it:} \emph{shishu}
  \item \emph{Say it:} \emph{shishu}, once more
  \item \emph{Say it:} say \emph{shishu}
  \item \emph{Recall:} say the Bengali for a husband, then the Bengali for a wife, then say \emph{shishu} again
\end{itemize}

\begin{tcolorbox}[breakable,colback=teal!4,colframe=teal!35!black,title={Before you move on}]
What does \bn{শিশু} mean? (\textquotedblleft{}A child\textquotedblright{}.) Say it once more.
\end{tcolorbox}
\section[\texorpdfstring{\bn{লোক} (lok)}{লোক (lok)}]{\bn{লোক} (lok) — a person}

\label{lesson:BN-C47-lok}



\begin{quote}
Before the new one: say the Bengali for a wife, then the Bengali for a child.
\end{quote}

\subsection*{You'll want to know: \bn{লোক}}
\textbf{\bn{লোক}} — \emph{lok} — \textquotedblleft{}a person\textquotedblright{}.

\begin{cousinweb}
\textbf{\bn{লোক}} (\emph{lok}) is Sanskrit \emph{loka}, \textquotedblleft{}world, people\textquotedblright{}, taken over whole as a learned loan, a tatsama (CDIAL 11119). The word inherited from it is Old Bengali \emph{loa}, which CDIAL lists.
\end{cousinweb}

\begin{grammarlens}[title={What you've built}]
One more word for this chapter.
\end{grammarlens}

\subsection*{Your turn}
\begin{itemize}
\raggedright
  \item \emph{Say it:} \emph{lok}
  \item \emph{Say it:} \emph{lok}, once more
  \item \emph{Say it:} say \emph{lok}
  \item \emph{Recall:} say the Bengali for a wife, then the Bengali for a child, then say \emph{lok} again
\end{itemize}

\begin{tcolorbox}[breakable,colback=teal!4,colframe=teal!35!black,title={Before you move on}]
What does \bn{লোক} mean? (\textquotedblleft{}A person\textquotedblright{}.) Say it once more.
\end{tcolorbox}
